To test whether the claims-calibrated sector adjustment (Section~\ref{ssec: sector}) transfers to real corporate exposures, we assemble an independent validation set of 24 flood events at publicly traded companies spanning six distinct floods, five countries, and nine sectors. No company in this set was used to calibrate any model parameter; the validation is fully out-of-sample.

\subsection{Validation sample}

Companies are selected on three criteria: (i)~physical facilities in a documented flood zone; (ii)~publicly available loss disclosures in SEC filings (10-K, 10-Q, 8-K), earnings calls, or regulatory reports; and (iii)~identifiable building characteristics (number of stories, approximate replacement value) for $\kappa$ lookup. Flood depths and durations are assigned from FEMA flood extent reports, hydrological literature, and post-event surveys. Loss figures were verified against SEC EDGAR filings and, for Japanese companies, against TSE earnings releases denominated in yen and converted at prevailing exchange rates.

The 24 events cover:
\begin{itemize}
    \item \textbf{Hurricane Harvey} (2017, $\tau \approx 288$\,h): 8 companies across petrochemical, power generation, hospital, and light-industrial sectors in the Houston metro area.
    \item \textbf{Thailand floods} (2011, $\tau \approx 1000$--$1200$\,h): 8 Japanese electronics and automotive manufacturers with single-plant exposures in the Ayutthaya and Pathum Thani industrial estates (depths 1.5--3.4\,m).
    \item \textbf{Other events}: Hurricane Sandy (2012, Con Edison), Hurricane Helene (2024, Baxter International), Hurricane Ida (2021, Phillips~66 physical damage component), KZN mudslide (2022, Toyota~SA), Michigan flash flood (2022, Abbott Labs), Hurricane Florence (2018, International Paper), Storm Desmond (2015, United Biscuits).
\end{itemize}

\subsection{Methodology}

For each company, we: (1)~assign building characteristics (stories, coverage band) based on facility descriptions in annual reports and satellite imagery; (2)~look up $\kappa_{\text{claims}}$ from Table~\ref{tab: kappa claims}; (3)~apply the coverage-gradient factor $\rho_{\text{gradient}} = 0.25$ (Eq.~\ref{eq: rho gradient}) to obtain $\kappa_{\text{total}}$; (4)~compute the duration-dependent damage ratio $\DF(h,\tau)$ using the calibrated Richards function (Eq.~\ref{eq: duration richards}, Table~\ref{tab: calibrated params}); and (5)~estimate $\ECL$ via Eq.~\ref{eq: total ecl}.

The predicted loss is compared to the actual disclosed loss. We report the ratio (predicted / actual) for each company.

\subsection{Results}

\begin{table}[!htbp]
\centering
\small
\caption{Out-of-sample validation on 24 corporate flood events. $\kappa_{\text{cl}}$: claims-based sector adjustment from Table~\ref{tab: kappa claims}. $\kappa_{\text{tot}} = \kappa_{\text{cl}} \times \rho_{\text{gradient}}$. No corporate disclosures were used in calibration.}
\label{tab: oos results}
\resizebox{\textwidth}{!}{
\begin{tabular}{llccrrr}
\toprule
Company & Event & Fl. & $\kappa_{\text{tot}}$ & ECL (\$M) & Actual (\$M) & Ratio \\
\midrule
\multicolumn{7}{l}{\textit{Hurricane Harvey (2017, $\tau \approx 288$\,h)}} \\
NRG Energy           & Harvey   & 1 & 0.21 & 133 &  75 & 1.8$\times$ \\
LyondellBasell       & Harvey   & 1 & 0.21 & 246 & 200 & 1.2$\times$ \\
ExxonMobil           & Harvey   & 1 & 0.21 & 287 & 170 & 1.7$\times$ \\
CenterPoint Energy   & Harvey   & 1 & 0.21 & 144 & 125 & 1.2$\times$ \\
Olin Corp            & Harvey   & 1 & 0.21 &  39 &  43 & 0.9$\times$ \\
Huntsman Corp        & Harvey   & 1 & 0.21 &  64 &  50 & 1.3$\times$ \\
Phillips 66          & Harvey   & 1 & 0.21 &  54 &  60 & 0.9$\times$ \\
HCA Healthcare       & Harvey   & 3 & 0.14 & 107 &  50 & 2.1$\times$ \\
\midrule
\multicolumn{7}{l}{\textit{Thailand floods (2011, $\tau \approx 1000$--$1200$\,h)}} \\
Western Digital      & Thailand & 1 & 0.21 & 284 & 275 & 1.0$\times$ \\
Canon                & Thailand & 1 & 0.21 & 177 & 260 & 0.7$\times$ \\
Toshiba              & Thailand & 1 & 0.21 & 190 & 260 & 0.7$\times$ \\
Panasonic            & Thailand & 1 & 0.21 & 127 & 170 & 0.7$\times$ \\
Nikon                & Thailand & 2 & 0.16 &  93 & 140 & 0.7$\times$ \\
Minebea              & Thailand & 2 & 0.16 &  94 & 120 & 0.8$\times$ \\
Pioneer              & Thailand & 1 & 0.21 &  58 & 100 & 0.6$\times$ \\
Rohm                 & Thailand & 1 & 0.21 &  51 &  67 & 0.8$\times$ \\
Honda$^{*}$          & Thailand & 1 & 0.21 & 190 & 500 & 0.38$\times$ \\
\midrule
\multicolumn{7}{l}{\textit{Other events}} \\
Toyota SA            & KZN 2022      & 1 & 0.21 & 308 & 361 & 0.9$\times$ \\
Baxter Int'l         & Helene 2024   & 2 & 0.16 & 125 & 200 & 0.6$\times$ \\
Abbott Labs          & Michigan 2022 & 1 & 0.21 &  23 &  15 & 1.5$\times$ \\
Con Edison           & Sandy 2012    & 1 & 0.21 & 489 & 460 & 1.1$\times$ \\
Int'l Paper          & Florence 2018 & 1 & 0.21 &  70 &  35 & 2.0$\times$ \\
Phillips 66          & Ida 2021      & 1 & 0.21 & 240 & 400 & 0.6$\times$ \\
United Biscuits      & Desmond 2015  & 2 & 0.16 &  39 &  30 & 1.3$\times$ \\
\bottomrule
\end{tabular}
}
\begin{flushleft}
\footnotesize
$^{*}$Honda's disclosed loss includes substantial supply-chain propagation from hundreds of Thai-based suppliers, which is outside the scope of the building-level damage model.
\end{flushleft}
\end{table}

Of the 24 predictions, 23 fall within the 0.4--2.5$\times$ range (96\%), and 21 within 0.5--2.0$\times$ (88\%). The median ratio is $0.90\times$ and the mean $|\!\log(\text{ratio})|$ is 0.37, indicating approximate unbiasedness with moderate dispersion.

\subsection{Discussion}

Three observations emerge from the out-of-sample validation.

\textit{First, the claims-based $\kappa$ captures building-level variation without corporate data.} Multi-storey buildings receive systematically lower $\kappa$ values than single-storey buildings, and this discrimination is reflected in the validation: the four two-storey companies (Baxter, Nikon, Minebea, United Biscuits) have a median ratio of $0.88\times$ with tight dispersion, compared to $1.14\times$ for single-storey companies. Because the storey count is derived from NFIP claims and not from the validation companies, this is a genuine out-of-sample signal.

\textit{Second, the coverage-gradient extrapolation provides a physically plausible resilience adjustment.} The estimated $\rho_{\text{gradient}} = 0.25$ implies that large corporate facilities absorb approximately 75\% of the damage that an NFIP-insured commercial building of the same height would sustain. This is consistent with the engineering intuition that large operators benefit from higher floor elevations, reinforced construction, comprehensive property insurance, and geographic redundancy. The resulting $\kappa_{\text{total}}$ for single-storey industrial facilities ($\approx 0.21$) falls between the literature range of 0.4--0.7 for SMEs~\cite{merz2010review, huizinga2017global} and the lower values expected for Fortune 500 companies.

\textit{Third, the model's scope limitations are transparent.} Honda ($0.38\times$) is the only prediction outside the 0.4--2.5$\times$ range: its disclosed \$500\,M loss includes supply-chain disruption from hundreds of Thai-based component suppliers, which the building-level damage model does not capture. HCA Healthcare ($2.1\times$) is the most overpredicted case, likely reflecting the model's difficulty with hospital closures where revenue loss and facility replacement interact in ways not captured by the standard CapEx/OpEx decomposition.

\paragraph{Geographic transferability} Although the calibration uses only US claims, the validation set spans five countries and three flood archetypes (tropical-cyclone surge in the US Gulf and Atlantic coasts; monsoonal/industrial-park flooding in Thailand; subtropical convective storm in KZN, South Africa; and mid-latitude winter storm in the UK). The two regional sub-samples large enough for a meaningful comparison (US Harvey, $n = 8$; Thailand, $n = 8$) yield prediction medians of $1.2\times$ and $0.78\times$ respectively, both well within the 0.4--2.5$\times$ acceptance band; the remaining events are single observations that we report as case studies rather than statistical evidence. This indirectly supports the coverage-gradient extrapolation against geographies and construction stocks not represented in the NFIP calibration sample, while we caution that per-country sample sizes outside the two main groups are too small to detect regional bias. The sample also remains weighted toward heavy-industrial and electronics manufacturing exposures; transferability to lighter commercial portfolios (retail, hospitality) and to non-OECD construction practices warrants targeted future testing.

\paragraph{Excluded events} Four events identified during data collection were excluded: (i)~Arkema Crosby (Harvey): chemical explosion triggered by refrigeration failure, not direct flood damage; (ii)~Kroger Houston (Harvey): losses dominated by perishable-inventory write-offs; (iii)~Entergy Louisiana (Ida): \$2.15\,B reflects transmission-grid rebuild, not building damage; (iv)~Phillips~66 Alliance impairment (Ida): the \$1.3\,B total includes a book-value write-down from permanent closure; only the estimated physical-damage component (\$400\,M) is retained. The exclusion criteria were specified before the predicted/actual ratios were computed (each excluded event has a documented loss-mechanism mismatch with the building-level damage model rather than a statistical outlier status), but we acknowledge that any post-hoc filter introduces selection-bias risk; the headline accuracy figures should be read with this caveat.
